%-------------------------
% Long CV -- Italian. Layout and shared macros live in preamble.tex.
% Keep resume.tex in sync with every content change made here.
%-------------------------
\documentclass[letterpaper,11pt,italian]{article}
%-------------------------
% Resume in Latex
% Author : Sourabh Bajaj
% Website: https://github.com/sb2nov/resume
% License : MIT
%------------------------
%
% Shared preamble for every variant of the long CV.
%
%   resume.tex            English
%   resume-it.tex         Italian
%   resume-signed.tex     English + GDPR clause and signature
%   resume-it-signed.tex  Italian + GDPR clause and signature
%
% The document language is chosen by the class option passed in
% \documentclass (english / italian) and picked up by babel.
% A signed variant defines \cvsigned before loading the CV.
%------------------------

\usepackage{babel}
\usepackage{fontawesome5}
\usepackage{latexsym}
\usepackage[empty]{fullpage}
\usepackage{titlesec}
\usepackage{marvosym}
\usepackage[usenames,dvipsnames]{color}
\usepackage{verbatim}
\usepackage{enumitem}
\usepackage{graphicx}
\usepackage[hidelinks]{hyperref}
\usepackage{fancyhdr}
\usepackage{comment}

\hypersetup{
  pdfauthor={Alessandro Monticelli},
  pdfsubject={Curriculum Vitae},
  pdfkeywords={Alessandro Monticelli, CV, Software Engineer, Embedded Systems, LoRa, RAG},
}

\pagestyle{fancy}
\fancyhf{} % clear all header and footer fields
\fancyfoot{}
\renewcommand{\headrulewidth}{0pt}
\renewcommand{\footrulewidth}{0pt}

% Adjust margins
\addtolength{\oddsidemargin}{-0.375in}
\addtolength{\evensidemargin}{-0.375in}
\addtolength{\textwidth}{1in}
\addtolength{\topmargin}{-.6in}
\addtolength{\textheight}{1.3in}

\urlstyle{same}

\raggedbottom
\raggedright
\setlength{\tabcolsep}{0in}

% Sections formatting
\titleformat{\section}{
  \scshape\raggedright\large
}{}{0em}{}[\color{black}\titlerule \vspace{4pt}]
\titlespacing*{\section}{0pt}{4pt}{3pt}

%-------------------------
% Custom commands
\newcommand{\resumeItem}[2]{%
  \item\small{%
    \textbf{#1}{: #2}%
  }%
}

\newcommand{\resumeSubheading}[4]{%
  \vspace{1pt}\item%
    \begin{tabular*}{0.97\textwidth}{l@{\extracolsep{\fill}}r}%
      \textbf{#1} & #2 \\%
      \textit{\small#3} & \textit{\small #4} \\%
    \end{tabular*}%
}

\newcommand{\resumeSubItem}[2]{\resumeItem{#1}{#2}}

\renewcommand{\labelitemii}{$\circ$}

\newcommand{\resumeSubHeadingListStart}{\begin{itemize}[leftmargin=*,noitemsep,topsep=1pt]}
\newcommand{\resumeSubHeadingListEnd}{\end{itemize}}
\newcommand{\resumeItemListStart}{\begin{itemize}[leftmargin=12pt,noitemsep,topsep=1pt]}
\newcommand{\resumeItemListEnd}{\end{itemize}}

%-------------------------
% Signed variants: GDPR consent clause, place/date and signature.
%
% \cvsignatureblock prints nothing unless \cvsigned is defined
% (see resume-signed.tex / resume-it-signed.tex).
%
% The signature image is read from signature.png when it exists.
% In CI (.github/workflows/signed.yml) it is decrypted from the committed
% signature.png.gpg with the SIGNATURE_PASSPHRASE secret, so the plain
% scan is never committed. Without it a blank signature line is printed.
%
% Set the place shown before the date with \def\cvplace{City}
% (left empty, only the date is printed).
%-------------------------

\ifdefined\cvplace\else\def\cvplace{}\fi
% Vertical space before the block and height of the signature image;
% a preamble can \def these before %-------------------------
% Signed variants: GDPR consent clause, place/date and signature.
%
% \cvsignatureblock prints nothing unless \cvsigned is defined
% (see resume-signed.tex / resume-it-signed.tex).
%
% The signature image is read from signature.png when it exists.
% In CI (.github/workflows/signed.yml) it is decrypted from the committed
% signature.png.gpg with the SIGNATURE_PASSPHRASE secret, so the plain
% scan is never committed. Without it a blank signature line is printed.
%
% Set the place shown before the date with \def\cvplace{City}
% (left empty, only the date is printed).
%-------------------------

\ifdefined\cvplace\else\def\cvplace{}\fi
% Vertical space before the block and height of the signature image;
% a preamble can \def these before \input{signature} to save space.
\ifdefined\cvsignatureskip\else\def\cvsignatureskip{8pt}\fi
\ifdefined\cvsignatureheight\else\def\cvsignatureheight{1cm}\fi

\newcommand{\cvconsenttext}{%
  \iflanguage{italian}{%
    Autorizzo il trattamento dei dati personali contenuti nel presente CV
    ai sensi del D.Lgs.~196/2003 e del Regolamento (UE) 2016/679 (GDPR).%
  }{%
    I authorize the processing of the personal data contained in this CV
    pursuant to Legislative Decree 196/2003 and EU Regulation 2016/679 (GDPR).%
  }%
}

\newcommand{\cvsignatureblock}{%
  \ifdefined\cvsigned
    \par\vspace{\cvsignatureskip}%
    \noindent\begin{minipage}[b]{0.66\linewidth}%
      \footnotesize\raggedright\cvconsenttext\par\vspace{2pt}%
      \ifx\cvplace\empty\else\cvplace, \fi\today
    \end{minipage}\hfill
    \begin{minipage}[b]{0.3\linewidth}%
      \centering
      \IfFileExists{signature.png}%
        {\includegraphics[height=\cvsignatureheight,width=\linewidth,keepaspectratio]{signature.png}}%
        {\rule{0pt}{\cvsignatureheight}}\\[-2pt]
      \rule{\linewidth}{0.4pt}\\
      \footnotesize Alessandro Monticelli
    \end{minipage}%
  \fi
}
 to save space.
\ifdefined\cvsignatureskip\else\def\cvsignatureskip{8pt}\fi
\ifdefined\cvsignatureheight\else\def\cvsignatureheight{1cm}\fi

\newcommand{\cvconsenttext}{%
  \iflanguage{italian}{%
    Autorizzo il trattamento dei dati personali contenuti nel presente CV
    ai sensi del D.Lgs.~196/2003 e del Regolamento (UE) 2016/679 (GDPR).%
  }{%
    I authorize the processing of the personal data contained in this CV
    pursuant to Legislative Decree 196/2003 and EU Regulation 2016/679 (GDPR).%
  }%
}

\newcommand{\cvsignatureblock}{%
  \ifdefined\cvsigned
    \par\vspace{\cvsignatureskip}%
    \noindent\begin{minipage}[b]{0.66\linewidth}%
      \footnotesize\raggedright\cvconsenttext\par\vspace{2pt}%
      \ifx\cvplace\empty\else\cvplace, \fi\today
    \end{minipage}\hfill
    \begin{minipage}[b]{0.3\linewidth}%
      \centering
      \IfFileExists{signature.png}%
        {\includegraphics[height=\cvsignatureheight,width=\linewidth,keepaspectratio]{signature.png}}%
        {\rule{0pt}{\cvsignatureheight}}\\[-2pt]
      \rule{\linewidth}{0.4pt}\\
      \footnotesize Alessandro Monticelli
    \end{minipage}%
  \fi
}


\hypersetup{pdftitle={Alessandro Monticelli -- Curriculum Vitae (IT)}, pdflang={it-IT}}

\begin{document}
\setlength{\footskip}{5pt}
%----------HEADING-----------------
\centering\textbf{\Large Alessandro Monticelli} \\
\vspace{0.08in}
\centering{\small{\faPhone} +39-340-053-1380 \quad
  \href{mailto:alessandro.monticelli@live.it}{\faEnvelope\ \underline{alessandro.monticelli@live.it}} \quad
  \href{https://www.linkedin.com/in/ale-mont}{{\faLinkedin} \underline{Alessandro Monticelli}} \quad
  \href{https://github.com/aleemont1}{{\faGithub} \underline{aleemont1}}}

%-----------EXPERIENCE-----------------
\section{Esperienze professionali}
\resumeSubHeadingListStart
\resumeSubheading
{Sviluppatore Software}{Faenza, Italia}
{S.EL.E.D. Soc. Coop.}{Settembre 2026 -- In corso}
\resumeItemListStart
  \resumeItem{Tecnologie utilizzate}{Python, FastAPI, PostgreSQL (pgvector, tsvector), Claude API (Prompt Caching), scikit-learn, Power Automate, Power Apps, SharePoint}
  \item\small{Progettato una piattaforma RAG ibrida a due stadi (FastAPI, PostgreSQL \texttt{pgvector} \& \texttt{tsvector}, RRF, Cross-Encoder) e un intent router basato su logit, con un modello di sicurezza air-gapped senza accesso diretto dell'LLM al database.}
  \item\small{Implementato query analitiche trasversali sui contratti e il monitoraggio automatico delle disdette, eliminando le scadenze di rinnovo mancate che causavano penali a cinque cifre per Caviro; disaccoppiato il motore principale per gestire anche i rapporti di prova di laboratorio.}
  \item\small{Ridotto di 15 volte i costi API di acquisizione documenti nelle pipeline ETL di Power Automate migrando a Claude Haiku 5.5, con prompt caching e output strutturati via tool calling al posto di lunghi prompt con schemi JSON.}
  \item\small{Rifattorizzato una PowerApp per la generazione degli XML delle fatture estere verso lo SdI (Agrintesa), con aggiornamenti massivi di aliquote IVA e codici \textit{Natura IVA}, validazione dei campi e UX ottimizzate.}
\resumeItemListEnd
\resumeSubheading
{Software Engineer freelance}{Da remoto}
{Libero professionista}{Giugno 2026 -- In corso}
\resumeItemListStart
  \resumeItem{Tecnologie utilizzate}{Python, SQLite, Docker, Multiprocessing, REST API}
  \item\small{Sviluppato \href{https://github.com/aleemont1/supermarket-scraper}{GDOScraper}, una pipeline ETL modulare che estrae, normalizza e salva i dati promozionali delle catene GDO italiane per la startup Foodify.}
  \item\small{Implementato lo scraping dinamico di REST API e il parsing visuale di volantini PDF, con segmentazione spaziale e ritaglio delle immagini tramite bounding box.}
  \item\small{Progettato un'architettura scalabile (Strategy Pattern, multiprocessing) con dashboard web e TUI per la verifica visuale dei dati.}
\resumeItemListEnd
\resumeSubheading
{Studente ricercatore}{Cesena, Italia}
{Sotto la supervisione del prof. Andrea Piroddi, Università di Bologna}{Luglio 2025 -- In corso}
\resumeItemListStart
\resumeItem{Tecnologie utilizzate} {C++17, LoRa\textsuperscript{\textcopyright}, ns-3, Unity (TDD), Semtech SX1262}
\resumeItem{Descrizione}
{
  \begin{itemize}
    \item Progettato \href{https://github.com/aleemont1/LAMP}{LAMP (LoRaMultiPacket)}, protocollo di trasporto C++17 indipendente dall'hardware, con header logico di 9 byte e SACK-ARQ a singola bitmap cumulativa, per payload fino a 62,2~KB.
    \item Condotto una campagna di simulazione stocastica in ns-3 (295.800 esecuzioni su canali statici, urbani e rurali), estendendo del 22--35\% la portata al 90\% di PDR in ambito urbano e aumentando il goodput del 14--16\% a SF7.
    \item Validato su banco hardware ESP32-S3 / SX1262 a doppia scheda con scostamenti entro il 7\%; articolo IEEE in stesura con il prof. Andrea Piroddi.
  \end{itemize}
}
\resumeItemListEnd
\resumeSubheading{Project Manager}{Forlì, Italia}{Aurora Rocketry Team - Unibo}{Aprile 2024 -- In corso}
\resumeItemListStart
\resumeItem{Tecnologie utilizzate} {C++, LoRa\textsuperscript{\textcopyright}, Python}
\resumeItem{Descrizione}
{\begin{itemize}
    \item Guido lo sviluppo del flight computer ESP32 per \href{https://github.com/AuroraRocketryTeam/atlas}{Atlas / Nemesis} (razzo sonda), coordinando un team di 7 sviluppatori.
    \item Supervisionato l'integrazione della telemetria con un protocollo LoRa\textsuperscript{\textcopyright} personalizzato, la stazione di terra e gli algoritmi di volo.
    \item Il sistema è alla base della mia tesi triennale ed è stato validato con successo in test statici e simulazioni.
  \end{itemize}}
\resumeItemListEnd
\resumeSubheading{Tirocinante -- Embedded Software Engineer}{Cesena, Italia}{SPOT. Digital Made Simple}{Febbraio 2024 -- Maggio 2024}
\resumeItemListStart
  \resumeItem{Tecnologie utilizzate}{Python, .NET 8.0, NuGet}
  \item\small{Reverse engineering dei driver dei moduli RF di poltrone smart per cinema, con estrazione dei comandi di controllo a basso livello.}
  \item\small{Progettato e sviluppato \href{https://github.com/aleemont1/OKINsimulator}{OKIN Simulator}, che replica il sistema IoT e ha ridotto del 60\% i tempi dei test di integrazione.}
  \item\small{Contribuito al refactoring del backend della piattaforma IoT (Python, .NET 8.0).}
\resumeItemListEnd


\newpage
\resumeSubheading{Sviluppatore Software}{Bologna, Italia}{UBM - Unibo Motorsport Team}{Ottobre 2023 -- Aprile 2024}
\resumeItemListStart
  \resumeItem{Tecnologie utilizzate}{Python, Vue.js, MySQL}
  \item\small{Migliorato il connection pooling del database, aumentando la stabilità delle connessioni e le prestazioni delle query.}
  \item\small{Ottimizzato gli algoritmi di scansione QR per la collocazione dei prodotti in magazzino, migliorandone l'affidabilità.}
\resumeItemListEnd
\resumeSubHeadingListEnd

%-----------EDUCATION-----------------
\section{Istruzione}
\resumeSubHeadingListStart
\resumeSubheading
{``Alma Mater Studiorum'' - Università di Bologna}{Cesena (FC), Italia}
{Laurea Magistrale in Ingegneria e Scienze Informatiche}{Settembre 2025 -- In corso}
\resumeSubheading
{Reykjavik University (H\'ask\'olinn \'i Reykjav\'ik)}{Reykjavík, Islanda}
{Scambio Erasmus+ nell'ambito della Laurea Magistrale in Computer Science}{Gennaio 2026 -- Maggio 2026}
\resumeItemListStart
\resumeItem{Insegnamenti rilevanti}{Distributed and Concurrent Programming (Go), Cybersecurity Management and Compliance in Practice (ISO27001, ISO42001), Deep Learning, Big Data Management}
\resumeItemListEnd
\resumeSubheading
{``Alma Mater Studiorum'' - Università di Bologna}{Cesena (FC), Italia}
{Laurea Triennale in Ingegneria e Scienze Informatiche}{Settembre 2021 -- Luglio 2025}
\resumeItemListStart
\resumeItem{Tesi}{``Sviluppo e ottimizzazione di un sistema di telemetria per razzi - Architettura e trasmissione dati tramite LoRa\textsuperscript{\textcopyright}''}
\resumeItem{Insegnamenti rilevanti}{Sistemi Embedded e Internet of Things, Programmazione di Reti, Reti di Telecomunicazione, High Performance Computing (CUDA, OpenMP, MPI), Integrazione e Virtualizzazione dei Sistemi, Programmazione ad Oggetti (Java, C\#)}
\resumeItemListEnd
\resumeSubHeadingListEnd

%--------PROGRAMMING SKILLS------------
\section{Competenze}
\resumeSubHeadingListStart
\item{
            \textbf{Programmazione}{: C++, C, Python, Java, C\# (.NET 8.0), Go (base), SQL.}
            \hfill
      }
\item{
            \textbf{AI \& Data Engineering}{: Architetture RAG, PostgreSQL (pgvector, tsvector), FastEmbed, scikit-learn, API di LLM (Prompt Caching), REST API (FastAPI).}
            \hfill
      }
\item{
            \textbf{Sistemi Embedded \& IoT}{: ESP32 (FreeRTOS), Arduino, Raspberry Pi, LoRa\textsuperscript{\textcopyright}, MQTT, Sensor Fusion, scheduling RTOS.}
            \hfill
      }
\item{
            \textbf{HPC \& Simulazione}{: OpenMP, OpenMPI, CUDA, SIMD, ns-3.}
            \hfill
      }
\item{
            \textbf{Cloud \& Integrazione di Sistemi}{: Docker, Kubernetes (base), CI/CD (GitHub Actions), Power Platform, Linux (Desktop \& Server), Windows Server.}
            \hfill
      }
\resumeSubHeadingListEnd

%-----------PROJECTS-----------------
\section{Progetti}
\resumeSubHeadingListStart
\resumeSubItem{\href{https://github.com/AuroraRocketryTeam/atlas}{Atlas / Nemesis (Flight Computer)}}
{Software di volo ESP32 per i razzi sonda di Aurora Rocketry: telemetria LoRa\textsuperscript{\textcopyright} in tempo reale, sensor fusion (IMU, barometro, GPS), concorrenza FreeRTOS e rilevamento dell'apogeo con EKF.}
\resumeSubItem{\href{https://github.com/aleemont1/federated-lightweight-chat}{Federated Lightweight Chat (FLC)}}
{Messaggistica distribuita peer-to-peer con vector clock, ordinamento causale, anti-entropia gossip, repliche SQLite e orchestrazione Docker con Traefik e Let's Encrypt.}
\resumeSubItem{\href{https://github.com/aleemont1/TheCouncil}{TheCouncil}}
{Framework open source di agentic engineering che offre code review in forma di tribunale dialettico e ``lenti'' di vincoli architetturali (Carmack, Lamport, Muratori, Ousterhout) per gli assistenti di programmazione AI.}
\resumeSubItem{\href{https://github.com/progetto-oop-22-23/OOP22-HOMER}{HOMER}}
{Progettato l'architettura OOP modulare e sviluppato gli algoritmi di monitoraggio energetico e automazione in tempo reale di una piattaforma smart home in Java.}
\resumeSubHeadingListEnd

%--------PUBLICATIONS--------
\section{Pubblicazioni}
\resumeSubHeadingListStart
\resumeSubItem{\href{https://github.com/aleemont1/lora-paper}{Articolo su rivista (in preparazione)}}
{Monticelli A., Piroddi A., \textit{“LAMP: A Selective-ACK Transport Protocol for Reliable Bulk Data Transmission over LoRa Networks”}, in preparazione per la sottomissione a IEEE, 2026.}
\resumeSubItem{\href{https://aleemont1.github.io/bachelor-thesis/pdf/tesi.pdf}{Tesi di Laurea Triennale}}
{\textit{“Sviluppo e ottimizzazione di un sistema di telemetria per razzi – Architettura e trasmissione dati tramite LoRa\textsuperscript{\textcopyright}”}, Università di Bologna, 2025.}
\resumeSubHeadingListEnd

%--------CERTIFICATES--------
\section{Certificazioni}
\resumeSubHeadingListStart
\item \textbf{Cisco Networking Academy}{: \href{https://github.com/aleemont1/curriculum-vitae/blob/main/certificates/PYTHON-2020_certificate.pdf}{PCAP -- Programming Essentials in Python}, \href{https://github.com/aleemont1/curriculum-vitae/blob/main/certificates/IOT-2018_certificate.pdf}{Introduction to the Internet of Things}, \href{https://github.com/aleemont1/curriculum-vitae/blob/main/certificates/GetConnected-2017_certificate.pdf}{Get Connected}.}
\item \textbf{Inglese}{: \href{https://github.com/aleemont1/curriculum-vitae/blob/main/certificates/English_C1.pdf}{Certificato di livello C1}.}
\resumeSubHeadingListEnd
\cvsignatureblock
%-------------------------------------------
\end{document}
